%% SECTION 14: LIMITS AND CONCLUSION
%% ============================================================
\section{Limits and Conclusion}

This paper does not prove the full economic thesis. It does not prove behavior preservation at corpus scale, does not prove paired training-token reduction, does not prove that a 100B-class canonical specialist replaces all broad frontier coding systems, does not apply to all novel systems or unsupported domains, and does not prove that behavior cells cover 70\%--90\% of routine product behavior. The QLoRA pilot is deliberately narrower: canonical translated trajectories are learnable under a parameter-efficient adaptation setup, and the measured forbidden-language markers remain at zero. The rest of the thesis must be earned through the claim ledger in Table~\ref{tab:claim-status}, beginning with the same-model raw/canonical substrate test.

The research agenda is therefore concrete. Build paired raw/canonical repositories from the same lineage. Assign every source artifact a legal/provenance disposition. Port behavior only within declared contracts. Report preservation tiers, weak-oracle failures, compatibility ghosts, and accepted incompatibilities. Run the same broad model on raw and canonical tasks before training a specialist. Measure files opened, context tokens, reasoning tokens, tool calls, invalid edits, failed proof lanes, reviewer comments, wall-clock, and cost per accepted proof-carrying change. Then run paired scaling curves before claiming training-token compression.

The thesis remains bold because the target is not cleaner code. The target is a canonical mirror of the software universe: behavior preserved or dispositioned, routine product work collapsed into certified cells, changes expressed as typed operations, proof lanes generated and checked, provenance carried through every transformation, and runtime failures compiled into negative memory. Human code contains the behavior, edge cases, incidents, product judgment, and hard-won lessons that matter. Raw human representation is the wrong thing to imitate when it encodes local accident rather than durable behavior.

The theoretical limit is the No-Accident Horizon developed in Section~\ref{sec:no-accident-horizon}. In that limit, agents do not invent local architecture, rediscover repository folklore, guess migration law, or re-learn the same auth, billing, upload, policy, and audit patterns in every codebase. They operate a constrained software machine and spend compute only where novelty, evidence, governance, risk, and future optionality remain.

The end state is no accidental software: no accidental architecture, dependency choice, CI ritual, migration path, security policy, review burden, or repeated reasoning once the substrate has learned the law. This paper is a falsifiable program for replacing accidental software representation with a canonical, proof-carrying substrate for verified correct change. In the strongest form, every remaining unit of engineering buys novelty, judgment, evidence, governance, or safety.

%% ============================================================
